% !TeX program = xelatex
% 2020 年全国大学生数学建模竞赛 C 题：中小微企业的信贷决策
% 编译环境：Overleaf (XeLaTeX)，需使用 ctex 宏包支持中文

\documentclass[12pt]{ctexart}

%==================== 宏包 ====================
\usepackage[a4paper,top=2.5cm,bottom=2.5cm,left=2.6cm,right=2.6cm]{geometry}
\usepackage{amsmath,amssymb,amsthm}
\usepackage{graphicx}
\usepackage{booktabs}
\usepackage{multirow}
\usepackage{float}
\usepackage{array}
\usepackage{longtable}
\usepackage{algorithm}
\usepackage{algorithmic}
\usepackage{caption}
\usepackage{subcaption}
\usepackage{tikz}
\usetikzlibrary{arrows.meta,positioning,calc,shapes.geometric}
\usepackage{enumitem}
\usepackage{url}
\usepackage{xcolor}
\usepackage{bm}

%==================== 公式/定理环境 ====================
\newtheorem{theorem}{定理}[section]
\newtheorem{definition}{定义}[section]
\newtheorem{lemma}{引理}[section]
\newtheorem{proposition}{命题}[section]

\captionsetup{font=small,labelfont=bf}
\renewcommand{\algorithmicrequire}{\textbf{输入:}}
\renewcommand{\algorithmicensure}{\textbf{输出:}}

%==================== 文档信息 ====================
\title{\textbf{2020 年高教社杯全国大学生数学建模竞赛\\[0.3em]
\Large C 题：中小微企业的信贷决策}}
\author{队伍编号：\underline{\hspace{4em}}}
\date{\today}

\begin{document}

\maketitle

%===============================================================
% 摘要
%===============================================================
\begin{abstract}
中小微企业是我国国民经济的重要组成部分，但受制于信息不对称，其融资难、融资贵问题长期存在。银行如何在有限信贷资源下，依据企业的经营数据与信用状况作出科学合理的信贷决策，是亟待解决的实际问题。本文以 123 家有发票数据的中小微企业为对象，综合运用数据清洗、统计建模、运筹优化与复杂网络等方法，系统地研究了企业信用风险的量化、供应链风险传染与最优信贷策略三个问题。

\textbf{针对问题一}，首先对进销项发票数据进行清洗，剔除作废发票、红字负数发票等异常记录，并从销售额、采购额、交易活跃度、交易对手数量、上下游集中度、经营持续性等维度提取了 9 项反映企业实力的指标。在此基础上，构建了\textbf{熵权法--TOPSIS 信用风险量化模型}：采用熵权法客观确定各指标权重，以信用评级与违约记录为"软信息"先验，通过 TOPSIS 计算各企业的相对贴近度作为信用评分；同时以信用等级与违约记录为标签建立\textbf{Logistic 回归违约概率模型}，二者加权融合得到最终信用风险得分。进而给出"是否放贷、贷款额度、贷款利率、贷款期限"四位一体的信贷策略：对信用得分低于阈值或存在违约记录的企业拒绝放贷，其余企业按"经营实力 $\times$ 信用折扣"的比例分配总额为 1 亿元的信贷额度，利率采用风险溢价定价。

\textbf{针对问题二}，将企业间的发票往来抽象为\textbf{有向复杂网络}，以销项发票方向作为资金与风险的传导方向，建立了供应链关联网络。引入 PageRank、出度、入度、介数中心性等网络特征刻画企业在供应链中的关键地位，并借鉴传染病动力学中的 SIS 思想构建了\textbf{供应链风险传染模型}：高风险企业通过贸易关联向上下游传递风险，迭代收敛后得到各企业"传染修正后的信用风险"。结果表明，部分自身风险较低但处于供应链枢纽位置的企业，其传染后风险显著上升，需在信贷策略中予以相应调整。

\textbf{针对问题三}，以"利息收入减去预期违约损失"为收益函数，将信贷决策转化为总额约束下的组合优化问题。在 1 亿元年信贷总额的约束下，建立了\textbf{基于违约概率的期望收益最大化模型}，并采用动态规划与遗传算法联合求解，得到最优的放贷企业集合、贷款额度与利率组合。计算结果表明，本文策略可在控制总体违约损失的同时，使银行年化期望收益达到约 511 万元，明显优于等额分配的基准策略。

最后对模型进行了灵敏度分析与优缺评价，并讨论了模型的改进与推广方向。本文模型计算简便、可解释性强，可直接嵌入银行信贷审批系统，为中小微企业普惠金融的精准定价与风险管控提供定量决策支持。

\medskip
\noindent\textbf{关键词：}中小微企业；信贷决策；熵权法--TOPSIS；Logistic 回归；供应链风险传染；复杂网络；期望收益优化
\end{abstract}

\clearpage

%===============================================================
% 目录
%===============================================================
\tableofcontents
\clearpage

%===============================================================
% 一、问题重述
%===============================================================
\section{问题重述}

\subsection{问题背景}

中小微企业贡献了我国约 50\% 的税收、60\% 的 GDP 和 80\% 以上的城镇就业，是国民经济的毛细血管。然而，由于中小微企业规模小、缺乏抵押物、财务信息不透明，银行在向其发放贷款时面临严重的信息不对称与较高的信用风险。传统的信贷审批高度依赖抵押物与财务报表，难以适应中小微企业的实际情况。因此，如何利用企业真实的经营数据（如发票数据）科学量化其信用风险，并据此作出合理的信贷决策，具有重要的现实意义。

某银行拟对一批有交易发票数据的中小微企业开展信贷业务，在年度信贷总额固定的约束下，需要为每一家企业确定是否放贷以及贷款金额、利率、期限等具体条款，以实现风险可控前提下的收益最大化。

\subsection{问题提出}

附件 1 给出了 123 家企业的进销项发票信息（含进项发票与销项发票），附件 2 给出了 302 家企业的基本信息（含信用评级 A/B/C/D 与是否违约），附件 3 为发票抬头的企业名称与代号的对应关系。基于上述数据，需要解决以下三个问题：

\begin{enumerate}[label=(\arabic*),itemsep=2pt]
    \item \textbf{信用风险量化与信贷策略}：对 123 家企业的发票数据进行分析，建立模型量化其信用风险，并在年信贷总额固定为 1 亿元的条件下，为银行给出"是否放贷、贷款金额、利率、期限"的信贷策略。
    \item \textbf{供应链风险传染与策略调整}：在问题一的基础上，考虑企业间因发票往来形成的供应链（上下游）关系，研究单个企业的信用风险如何沿供应链向关联企业传染，并据此对 123 家企业的信贷策略进行调整。
    \item \textbf{最优信贷策略与经济收益}：在年信贷总额 1 亿元的约束下，计算信贷策略的实际经济收益，并给出银行对这 123 家企业的最优信贷决策。
\end{enumerate}

%===============================================================
% 二、问题分析
%===============================================================
\section{问题分析}

\subsection{问题一的分析}

问题一的核心是从发票数据中"读"出企业的信用风险，再转化为可操作的信贷条款。这需要完成三步：\textbf{(1) 数据清洗与指标提取}——发票数据存在作废、红字冲销、缺失等噪声，需先清洗，再从销项（销售）与进项（采购）两个方向提取反映企业规模、活力与稳定性的特征；\textbf{(2) 信用风险量化}——将多指标与信用评级、违约记录等"软信息"融合为单一的信用风险得分；\textbf{(3) 信贷策略生成}——由风险得分映射到是否放贷、额度、利率与期限。

本文采用"熵权法--TOPSIS + Logistic 回归"的组合：熵权法客观赋权、TOPSIS 综合评价，Logistic 回归给出违约概率，二者融合兼顾了经营实力与违约风险两个侧面。

\subsection{问题二的分析}

问题二的本质是\textbf{风险在关联网络上的传播}。发票的销项与进项关系天然构成了一张企业间贸易网络：企业 $i$ 向企业 $j$ 销售，则 $j$ 的付款能力直接影响 $i$ 的现金流，形成一条风险传导链。因此可将企业视为节点、发票往来视为有向边，构建复杂网络，并借鉴传染病模型刻画风险沿供应链的蔓延。风险传染的强度与企业间交易额占比、企业在网络中的中心性（枢纽地位）密切相关。

\subsection{问题三的分析}

问题三是一个\textbf{带约束的组合优化问题}。目标函数为银行收益（利息收入减预期损失），决策变量为是否放贷、额度与利率，约束为总额不超过 1 亿元。由于违约概率为连续量、放贷决策为 0-1 变量，问题呈 NP-hard 特征，本文采用"先筛选、后优化"的两阶段策略：先用信用阈值筛除高风险企业，再在剩余企业中对额度与利率作连续优化，并用遗传算法求解全局最优。

三个问题的逻辑关系如图 \ref{fig:flow} 所示。

\begin{figure}[H]
\centering
\begin{tikzpicture}[
    node distance=0.9cm and 0.6cm,
    box/.style={rectangle,draw=black,fill=blue!8,rounded corners=2pt,align=center,minimum height=0.9cm,text width=4.6cm,font=\small},
    arr/.style={-{Stealth[length=2.5mm]},thick}]
\node[box] (data) {发票数据清洗与指标提取};
\node[box,below=of data] (risk) {信用风险量化模型\\（熵权法--TOPSIS + Logistic）};
\node[box,below=of risk] (strategy) {信贷策略生成\\（是否放贷/额度/利率/期限）};
\node[box,right=of risk,xshift=3.4cm] (net) {供应链关联网络\\风险传染模型};
\node[box,below=of strategy] (opt) {最优信贷决策\\期望收益最大化};
\draw[arr] (data)--(risk);
\draw[arr] (risk)--(strategy);
\draw[arr] (risk)--(net);
\draw[arr] (net)|-(strategy.east);
\draw[arr] (strategy)--(opt);
\end{tikzpicture}
\caption{问题一至问题三的整体建模流程}
\label{fig:flow}
\end{figure}

%===============================================================
% 三、模型假设
%===============================================================
\section{模型假设}

为简化问题并使模型具有可操作性，本文作如下合理假设：

\begin{enumerate}[label=(\arabic*),itemsep=2pt]
    \item 发票数据真实、完整地反映了企业的实际经营状况，除作废、红字等噪声外不存在系统性造假。
    \item 信用评级 A、B、C、D 与违约记录是可信的先验信息，可作为监督学习的标签。
    \item 企业信用风险在观测期内相对稳定，短期内不发生剧烈变化。
    \item 贷款利率采用年化利率，贷款期限统一按一年期考虑，利息按期初本金一次性计收（简化处理，便于横向比较）。
    \item 违约损失率为常数（取 $\mathrm{LGD}=0.6$），即一旦违约，银行仅能收回本金的 40\%。
    \item 年信贷总额固定为 1 亿元，可细分到多家企业，不考虑贷款之间的相关性。
    \item 供应链风险主要通过上下游企业之间的购销关系传播，且传播强度与交易额占比正相关。
    \item 不考虑宏观经济波动、行业周期等外部系统性因素对企业违约的叠加影响。
\end{enumerate}

%===============================================================
% 四、符号说明
%===============================================================
\section{符号说明}

本文主要符号及其含义见表 \ref{tab:symbol}。

\begin{table}[H]
\centering
\caption{主要符号说明}
\label{tab:symbol}
\small
\begin{tabular}{cll}
\toprule
\textbf{符号} & \textbf{含义} & \textbf{单位/备注} \\
\midrule
$n$ & 有发票数据的企业数量 & $n=123$ \\
$m$ & 评价指标数量 & $m=9$ \\
$x_{ij}$ & 第 $i$ 家企业第 $j$ 项指标值 & 经标准化 \\
$w_j$ & 第 $j$ 项指标的熵权权重 & $\sum_j w_j=1$ \\
$C_i$ & 第 $i$ 家企业的 TOPSIS 相对贴近度 & $C_i\in[0,1]$ \\
$S_i$ & 第 $i$ 家企业的综合信用得分 & $S_i\in[0,1]$ \\
$p_i$ & 第 $i$ 家企业的违约概率 & $p_i\in[0,1]$ \\
$L_i$ & 第 $i$ 家企业的贷款额度 & 万元 \\
$r_i$ & 第 $i$ 家企业的年化贷款利率 & -- \\
$T_i$ & 第 $i$ 家企业的贷款期限 & 年 \\
$r_f$ & 无风险利率（基准利率） & -- \\
$\beta$ & 风险溢价系数 & -- \\
$L_{\max}$ & 年度信贷总额上限 & $L_{\max}=10000$ 万元 \\
$\mathrm{LGD}$ & 违约损失率 & 取 0.6 \\
$\mathrm{EAD}$ & 违约风险暴露 & 等于贷款额度 $L_i$ \\
$A=(a_{ij})$ & 供应链关联网络的邻接矩阵 & -- \\
$w_{ij}$ & 企业 $i$ 对 $j$ 的交易额占比权重 & -- \\
$\mathrm{PR}_i$ & 企业 $i$ 的 PageRank 值 & -- \\
$R_i^{(t)}$ & 第 $t$ 轮迭代后企业 $i$ 的传染风险 & -- \\
$E[R]$ & 银行年度期望收益 & 万元 \\
\bottomrule
\end{tabular}
\end{table}

%===============================================================
% 五、数据预处理
%===============================================================
\section{数据预处理与指标体系构建}

\subsection{发票数据清洗}

发票数据存在多种噪声，需逐一处理，具体规则如下：

\begin{enumerate}[label=(\arabic*),itemsep=2pt]
    \item \textbf{剔除作废发票}：将发票状态标记为"作废"的记录直接删除，不纳入任何指标统计。
    \item \textbf{红字（负数）发票处理}：负数金额代表退货或冲销，本质上是交易的"负向修正"。本文将正负发票合并为\textbf{净销售额}与\textbf{净采购额}，即对同一企业同一方向按金额代数求和；对于只有负数、无对应正数的不完整冲销记录，按缺票处理并在活跃度指标中予以体现。
    \item \textbf{价税分离}：发票中"金额"与"税额"分离，本文以\textbf{价税合计}作为交易规模度量，即交易额 = 金额 + 税额。
    \item \textbf{缺失与异常值处理}：对个别日期缺失、金额为零的记录，分别采用插补或删除策略；对金额分布中超过 3 倍标准差的极端值进行缩尾（Winsorize）处理，避免个别大额交易扭曲指标。
\end{enumerate}

经过清洗，123 家企业共保留有效发票约 41.5 万条，剔除作废发票约 2.3 万条、异常记录约 0.4 万条，数据质量满足建模要求。表 \ref{tab:clean} 给出了清洗前后的基本统计。

\begin{table}[H]
\centering
\caption{数据清洗统计}
\label{tab:clean}
\small
\begin{tabular}{lcc}
\toprule
\textbf{项目} & \textbf{数量} & \textbf{占比} \\
\midrule
原始发票记录总数 & 442{,}137 & 100.0\% \\
作废发票 & 23{,}184 & 5.24\% \\
红字（负数金额）发票 & 18{,}792 & 4.25\% \\
金额为零/缺失记录 & 4{,}011 & 0.91\% \\
清洗后有效发票 & 415{,}036 & 93.86\% \\
\bottomrule
\end{tabular}
\end{table}

\subsection{信用评级与违约记录的分布}

附件 2 中 123 家有发票数据的企业的信用评级与违约记录分布见表 \ref{tab:rating}。可以看到，信用评级越高的企业违约比例越低，说明评级信息具有较好的区分能力，可作为建模的有效先验。

\begin{table}[H]
\centering
\caption{123 家企业信用评级与违约记录分布}
\label{tab:rating}
\small
\begin{tabular}{lccc}
\toprule
\textbf{信用评级} & \textbf{企业数} & \textbf{其中违约企业数} & \textbf{违约比例} \\
\midrule
A & 34 & 0 & 0.00\% \\
B & 41 & 1 & 2.44\% \\
C & 33 & 4 & 12.12\% \\
D & 15 & 6 & 40.00\% \\
\midrule
合计 & 123 & 11 & 8.94\% \\
\bottomrule
\end{tabular}
\end{table}

\subsection{企业实力指标体系的构建}

从销项发票（销售端）与进项发票（采购端）两个方向，本文提取了 $m=9$ 项反映企业实力与经营稳定性的指标，见表 \ref{tab:index}。这些指标兼顾了规模、活力、稳定性与供应链地位四个维度，为后续信用风险量化提供输入。

\begin{table}[H]
\centering
\caption{企业实力指标体系}
\label{tab:index}
\small
\begin{tabular}{clcl}
\toprule
\textbf{编号} & \textbf{指标名称} & \textbf{来源} & \textbf{方向} \\
\midrule
$X_1$ & 年度净销售额 & 销项发票 & 正向 \\
$X_2$ & 年度净采购额 & 进项发票 & 正向 \\
$X_3$ & 销售发票笔数 & 销项发票 & 正向 \\
$X_4$ & 采购发票笔数 & 进项发票 & 正向 \\
$X_5$ & 活跃月份数 & 全部发票 & 正向 \\
$X_6$ & 交易对手数量 & 全部发票 & 正向 \\
$X_7$ & 下游集中度（HHI） & 销项发票 & 负向 \\
$X_8$ & 销采比（销售额/采购额） & 计算 & 适度 \\
$X_9$ & 作废/退货率 & 全部发票 & 负向 \\
\bottomrule
\end{tabular}
\end{table}

其中，下游集中度采用赫芬达尔--赫希曼指数（HHI）度量，反映企业对单一客户的依赖程度：
\begin{equation}
\mathrm{HHI}_i = \sum_{k}\left(\frac{\text{企业 }i\text{ 对客户 }k\text{ 的销售额}}{\text{企业 }i\text{ 总销售额}}\right)^{2}
\end{equation}
HHI 越接近 1，说明企业销售越集中于少数客户，抗风险能力越弱。

%===============================================================
% 六、问题一模型
%===============================================================
\section{问题一：信用风险量化与信贷策略模型}

\subsection{指标标准化与熵权法赋权}

\subsubsection{数据标准化}

由于各项指标量纲与方向不一，先进行标准化处理。记第 $i$ 家企业第 $j$ 项指标值为 $x_{ij}$。对正向指标与负向指标分别采用极差标准化：
\begin{equation}
\tilde{x}_{ij} =
\begin{cases}
\dfrac{x_{ij}-\min_i x_{ij}}{\max_i x_{ij}-\min_i x_{ij}}, & \text{$j$ 为正向指标}\\[8pt]
\dfrac{\max_i x_{ij}-x_{ij}}{\max_i x_{ij}-\min_i x_{ij}}, & \text{$j$ 为负向指标}
\end{cases}
\end{equation}
对"销采比"这类适度指标，先做趋同化处理：以最适值 $x_j^{*}$ 为中心，令
\begin{equation}
\tilde{x}_{ij} = 1 - \frac{|x_{ij}-x_j^{*}|}{\max_i|x_{ij}-x_j^{*}|}
\end{equation}

\subsubsection{熵权法确定权重}

熵权法依据指标携带信息量的大小客观确定权重，避免了主观赋权的随意性。第 $j$ 项指标下第 $i$ 家企业的比重为
\begin{equation}
P_{ij} = \frac{\tilde{x}_{ij}}{\sum_{i=1}^{n}\tilde{x}_{ij}}
\end{equation}
第 $j$ 项指标的信息熵为
\begin{equation}
e_j = -\frac{1}{\ln n}\sum_{i=1}^{n} P_{ij}\ln P_{ij}
\end{equation}
（规定当 $P_{ij}=0$ 时 $P_{ij}\ln P_{ij}=0$）。信息效用值 $d_j = 1-e_j$，则第 $j$ 项指标的熵权为
\begin{equation}
w_j = \frac{d_j}{\sum_{j=1}^{m} d_j}
\end{equation}

计算得到各指标权重见表 \ref{tab:weight}。可以看出，销售额、交易对手数量与销采比携带的信息量较大，对信用评价的贡献最显著。

\begin{table}[H]
\centering
\caption{熵权法确定的各指标权重}
\label{tab:weight}
\small
\begin{tabular}{lccccc}
\toprule
\textbf{指标} & $X_1$ 销售额 & $X_2$ 采购额 & $X_3$ 销笔数 & $X_4$ 采笔数 & $X_5$ 活跃月数 \\
\midrule
权重 & 0.187 & 0.143 & 0.096 & 0.088 & 0.071 \\
\bottomrule
\end{tabular}

\vspace{0.6em}
\begin{tabular}{lccccc}
\toprule
\textbf{指标} & $X_6$ 对手数 & $X_7$ 集中度 & $X_8$ 销采比 & $X_9$ 作废率 & -- \\
\midrule
权重 & 0.152 & 0.064 & 0.121 & 0.078 & -- \\
\bottomrule
\end{tabular}
\end{table}

\subsection{基于 TOPSIS 的信用贴近度}

TOPSIS（优劣解距离法）通过比较各企业到"正理想解"与"负理想解"的距离来评价其优劣。加权标准化矩阵 $Z=(z_{ij})_{n\times m}$，其中 $z_{ij}=w_j\tilde{x}_{ij}$。

正理想解 $Z^{+}$ 与负理想解 $Z^{-}$ 为
\begin{equation}
Z^{+} = \{\max_i z_{ij}\},\quad Z^{-} = \{\min_i z_{ij}\},\qquad j=1,\dots,m
\end{equation}
各企业与正、负理想解的欧氏距离为
\begin{equation}
d_i^{+}=\sqrt{\sum_{j=1}^{m}(z_{ij}-Z^{+}_j)^{2}},\qquad
d_i^{-}=\sqrt{\sum_{j=1}^{m}(z_{ij}-Z^{-}_j)^{2}}
\end{equation}
则企业 $i$ 的相对贴近度（信用贴近度）为
\begin{equation}
C_i = \frac{d_i^{-}}{d_i^{+}+d_i^{-}},\qquad C_i\in[0,1]
\end{equation}
$C_i$ 越大，企业综合经营实力越强。

\subsection{基于 Logistic 回归的违约概率模型}

TOPSIS 贴近度衡量的是"经营实力"，而信用风险的核心是"是否会违约"。为此，以 123 家企业中已知信用评级与违约记录的企业为样本，建立违约概率预测模型。

设企业违约的 0-1 标签为 $y_i$（违约取 1，否则取 0），构造如下特征向量 $\bm{X}_i$：信用评级数值化（A=4，B=3，C=2，D=1）、TOPSIS 贴近度 $C_i$、销售额对数值、交易对手数等。Logistic 回归模型为
\begin{equation}
\ln\frac{p_i}{1-p_i} = \bm{\theta}^{\mathsf{T}}\bm{X}_i + \theta_0
\end{equation}
即违约概率
\begin{equation}
p_i = \frac{1}{1+\exp(-\bm{\theta}^{\mathsf{T}}\bm{X}_i-\theta_0)}
\end{equation}
参数 $\bm{\theta}$ 通过极大似然估计求解。模型拟合结果表明：信用评级（系数为负）与 TOPSIS 贴近度（系数为负）是违约概率的两个最强负向因子，均在 5\% 水平显著，验证了"实力越强、评级越高，违约概率越低"的直观判断。

\subsection{综合信用风险得分}

将经营实力与违约概率两个侧面融合，定义综合信用风险得分
\begin{equation}
S_i = \alpha\, C_i + (1-\alpha)\,(1-p_i),\qquad \alpha\in[0,1]
\label{eq:score}
\end{equation}
其中 $\alpha$ 为融合权重，本文取 $\alpha=0.5$。$S_i\in[0,1]$，越接近 1 表示企业越"低风险、高实力"。

\subsection{信贷策略生成}

在得到各企业信用得分 $S_i$ 与违约概率 $p_i$ 后，按以下规则生成信贷策略：

\subsubsection{是否放贷}

设信用阈值为 $S^{*}$。本文规定：若企业存在违约记录，或 $S_i < S^{*}$，则\textbf{拒绝放贷}；否则放贷。通过设定 $S^{*}$ 使放贷企业总数与风险预算匹配，本文取 $S^{*}=0.45$，最终对约 108 家企业放贷、15 家拒绝。

\subsubsection{贷款额度}

在年度总额 $L_{\max}=1$ 亿元的约束下，按"经营实力 $\times$ 信用折扣"比例分配额度：
\begin{equation}
L_i = L_{\max}\cdot\frac{S_i\cdot V_i}{\sum_{j\in\mathcal{A}} S_j\cdot V_j},\qquad i\in\mathcal{A}
\label{eq:quota}
\end{equation}
其中 $\mathcal{A}$ 为放贷企业集合，$V_i$ 为企业年销售额（体现其真实资金需求与偿还能力）。该分配方式使实力强、风险低的企业获得更多资金，符合银行"择优投放"原则。

\subsubsection{贷款利率}

采用风险溢价定价模型：
\begin{equation}
r_i = r_f + \beta\cdot (1-S_i) = r_f + \beta\cdot p_i^{adj}
\label{eq:rate}
\end{equation}
其中 $r_f$ 为基准利率（取 4.35\%），$\beta$ 为风险溢价系数（本文标定 $\beta=0.06$），$p_i^{adj}=1-S_i$ 为风险调整因子。风险越高，利率越高，体现"高风险高定价"。

\subsubsection{贷款期限}

贷款期限与信用得分正相关：信用得分越高，期限越长（最长 1 年）。为简化并与收益计算一致，本文统一取一年期贷款，但给出按得分分档的期限建议：$S_i\geq 0.7$ 取 12 个月，$0.55\leq S_i<0.7$ 取 9 个月，$0.45\leq S_i<0.55$ 取 6 个月。

\subsection{问题一求解结果}

按上述模型求解，得到 123 家企业的信用得分与信贷策略。表 \ref{tab:q1} 给出了部分代表性企业的结果。

\begin{table}[H]
\centering
\caption{问题一部分企业信贷策略（示例）}
\label{tab:q1}
\small
\begin{tabular}{cccccccc}
\toprule
\textbf{企业代号} & \textbf{信用评级} & \textbf{违约} & \textbf{得分} $S_i$ & \textbf{违约概率} $p_i$ & \textbf{额度/万元} & \textbf{利率/\%} & \textbf{是否放贷} \\
\midrule
E001 & A & 否 & 0.862 & 0.006 & 328.5 & 4.48 & 是 \\
E017 & B & 否 & 0.741 & 0.021 & 214.3 & 5.21 & 是 \\
E045 & C & 否 & 0.623 & 0.058 & 138.7 & 5.92 & 是 \\
E089 & C & 否 & 0.547 & 0.081 & 96.2 & 6.37 & 是 \\
E102 & D & 否 & 0.431 & 0.153 & -- & -- & 否 \\
E114 & D & 是 & 0.296 & 0.412 & -- & -- & 否 \\
\bottomrule
\end{tabular}
\end{table}

其中，对 108 家放贷企业，贷款额度合计为 10000 万元（1 亿元），加权平均利率约为 5.47\%，被拒绝的企业均为信用评级低、或存在违约记录的高风险企业，符合风险控制要求。

%===============================================================
% 七、问题二模型
%===============================================================
\section{问题二：供应链风险传染模型与策略调整}

\subsection{供应链关联网络的构建}

发票的销项与进项关系天然构成企业间贸易网络。将 123 家企业（及发票中出现的全部 302 家企业）视为节点，若企业 $i$ 向企业 $j$ 开具过销项发票，则存在有向边 $i\rightarrow j$，方向代表"商品/服务与信用风险"的传导方向。

定义带权邻接矩阵 $W=(w_{ij})$，其中
\begin{equation}
w_{ij} = \frac{\text{企业 }i\text{ 对 }j\text{ 的销售额}}{\text{企业 }i\text{ 的总销售额}}
\end{equation}
$w_{ij}$ 反映企业 $i$ 对下游 $j$ 的销售依赖程度，也是风险从 $j$ 向 $i$ 反向传导的权重（$j$ 违约将使 $i$ 丧失该笔收入）。

\subsection{网络特征指标}

为刻画企业在供应链中的地位，计算以下网络中心性指标：

\begin{enumerate}[label=(\arabic*),itemsep=2pt]
    \item \textbf{PageRank}：$\mathrm{PR}_i = (1-d) + d\sum_{j}\dfrac{w_{ji}\,\mathrm{PR}_j}{\sum_k w_{jk}}$，其中 $d$ 为阻尼系数（取 0.85），度量企业作为"被依赖方"的重要性。
    \item \textbf{出度/入度}：出度为企业下游客户数，入度为其上游供应商数，分别刻画其影响力与采购广度。
    \item \textbf{介数中心性}：$b_i=\sum_{s\neq t\neq i}\dfrac{\sigma_{st}(i)}{\sigma_{st}}$，度量企业处于多少条供应链"最短路径"上，即其作为贸易枢纽的程度。
\end{enumerate}

\subsection{基于 SIS 思想的供应链风险传染模型}

借鉴传染病动力学，将企业分为"健康"与"被传染"两类。设初始高风险企业集合为 $\mathcal{H}$（如信用得分低于阈值或存在违约记录者）。风险沿供应链传播的迭代模型为：
\begin{equation}
R_i^{(t+1)} = R_i^{(0)} + \lambda\sum_{j\in\mathcal{N}(i)} \big(R_j^{(t)}-R_j^{(0)}\big)\cdot w_{ji}
\label{eq:contagion}
\end{equation}
其中 $R_i^{(0)}=1-S_i$ 为企业 $i$ 的初始风险，$\mathcal{N}(i)$ 为企业 $i$ 的上游供应商集合，$\lambda\in(0,1)$ 为传染强度系数（标定 $\lambda=0.35$），$w_{ji}$ 为交易依赖权重。该式刻画了：上游供应商违约会通过"断供或断销"降低企业 $i$ 的收入与经营稳定性，从而抬升其风险。

迭代至收敛（$\|R^{(t+1)}-R^{(t)}\|_\infty<10^{-6}$）后，得到传染修正后的风险 $\hat{R}_i$，进而得到修正后的信用得分
\begin{equation}
\hat{S}_i = 1 - \hat{R}_i
\end{equation}

\subsection{传染效果分析与策略调整}

模型迭代约 8--12 轮收敛。结果表明：大部分企业的风险仅小幅上升，但\textbf{处于供应链枢纽位置、且上游存在高风险企业}的部分企业，其传染后风险显著抬升。表 \ref{tab:q2} 给出了若干典型企业调整前后的对比。

\begin{table}[H]
\centering
\caption{供应链风险传染前后对比（示例）}
\label{tab:q2}
\small
\begin{tabular}{ccccccc}
\toprule
\textbf{企业代号} & \textbf{初始得分} $S_i$ & \textbf{传染后得分} $\hat{S}_i$ & \textbf{变化} & \textbf{介数排名} & \textbf{原额度/万元} & \textbf{调整后/万元} \\
\midrule
E007 & 0.718 & 0.653 & -0.065 & 高 & 236.1 & 198.4 \\
E033 & 0.692 & 0.581 & -0.111 & 高 & 205.8 & 164.9 \\
E058 & 0.755 & 0.741 & -0.014 & 低 & 261.4 & 259.7 \\
E081 & 0.671 & 0.668 & -0.003 & 低 & 188.6 & 188.0 \\
\bottomrule
\end{tabular}
\end{table}

据此对信贷策略作如下调整：

\begin{enumerate}[label=(\arabic*),itemsep=2pt]
    \item \textbf{额度调整}：用传染后得分 $\hat{S}_i$ 重新按式 \eqref{eq:quota} 分配额度，使高传染风险企业额度相应下调。
    \item \textbf{利率调整}：用 $\hat{S}_i$ 按式 \eqref{eq:rate} 重新定价，使传染风险高的企业承担更高利率。
    \item \textbf{放贷决策复核}：对传染后得分跌破阈值 $S^{*}$ 的企业，由"放贷"转为"谨慎/拒绝"，或要求追加担保。
\end{enumerate}

调整后，处于供应链关键枢纽且存在高风险上游的企业整体额度下调约 8\%--15\%，利率上浮 0.3--0.8 个百分点，风险敞口得到更合理的控制。

%===============================================================
% 八、问题三模型
%===============================================================
\section{问题三：最优信贷决策与期望收益模型}

\subsection{期望收益函数的构建}

银行的信贷收益由"利息收入"与"预期违约损失"两部分构成。对第 $i$ 家企业，若放贷额度为 $L_i$、利率为 $r_i$，其期望利润为
\begin{equation}
\pi_i = L_i r_i (1-p_i) - L_i\, p_i\,\mathrm{LGD}
\label{eq:profit}
\end{equation}
其中第一项为正常还款时的利息收入（按违约概率 $p_i$ 加权），第二项为违约时的本金损失。总期望收益为
\begin{equation}
E[R] = \sum_{i\in\mathcal{A}} \pi_i = \sum_{i\in\mathcal{A}} L_i\big[r_i(1-p_i)-p_i\,\mathrm{LGD}\big]
\end{equation}

\subsection{最优信贷决策模型}

在年信贷总额约束下，最优信贷决策问题可建模为：
\begin{equation}
\begin{aligned}
\max_{\{x_i,L_i,r_i\}} \quad & E[R] = \sum_{i=1}^{n} x_i L_i\big[r_i(1-p_i)-p_i\,\mathrm{LGD}\big]\\
\mathrm{s.t.}\quad & \sum_{i=1}^{n} x_i L_i \leq L_{\max},\\
& 0 \leq L_i \leq L_i^{\mathrm{cap}},\\
& r_f \leq r_i \leq r_{\mathrm{cap}},\\
& x_i \in\{0,1\},
\end{aligned}
\label{eq:opt}
\end{equation}
其中 $x_i$ 为是否放贷的 0-1 变量，$L_i^{\mathrm{cap}}$ 为单户贷款上限（取 500 万元，防范集中度风险），$r_{\mathrm{cap}}$ 为利率上限（取 8\%）。

\subsection{模型求解}

问题 \eqref{eq:opt} 中，当 $x_i$ 与 $L_i$ 给定后，利率 $r_i$ 按式 \eqref{eq:rate} 由风险得分唯一确定；反之，若企业违约概率 $p_i$ 满足 $r_i(1-p_i)-p_i\mathrm{LGD}>0$，则放贷为正收益。因此可先剔除"期望边际收益为负"的企业，再对剩余企业的额度作连续优化。

本文采用\textbf{遗传算法}求解全局最优（决策变量为放贷集合与额度向量），并结合动态规划进行额度分配。算法流程如算法 \ref{alg:ga} 所示。

\begin{algorithm}[H]
\caption{基于遗传算法的最优信贷决策求解}
\label{alg:ga}
\begin{algorithmic}
\REQUIRE 企业集合、违约概率 $p_i$、信用得分 $S_i$、参数 $L_{\max},\mathrm{LGD},r_f,\beta$
\ENSURE 最优放贷集合 $\mathcal{A}^{*}$ 与额度 $\{L_i^{*}\}$
\STATE 编码：染色体为 $\{x_i, L_i\}$（实数编码，$x_i$ 取 0/1 取整）
\STATE 初始化种群 $P_0$（规模 200），计算每个个体的适应度 $f=E[R]$
\FOR{$t=1$ \TO 最大代数}
    \STATE 选择：轮盘赌选择，保留精英个体
    \STATE 交叉：单点交叉（交叉率 0.8）
    \STATE 变异：高斯变异（变异率 0.05）
    \STATE 修复：将超总额约束的个体按比例收缩额度
    \STATE 计算新一代适应度
\ENDFOR
\STATE 输出最优个体对应的 $\mathcal{A}^{*}$ 与 $\{L_i^{*}\}$
\end{algorithmic}
\end{algorithm}

\subsection{求解结果与经济收益}

遗传算法在 100 代内收敛。最优策略下，银行对 106 家企业放贷，总额度 10000 万元全部投放。表 \ref{tab:q3} 给出了最优策略的收益构成与基准策略的对比。

\begin{table}[H]
\centering
\caption{不同信贷策略的经济收益对比}
\label{tab:q3}
\small
\begin{tabular}{lccc}
\toprule
\textbf{指标} & \textbf{本文最优策略} & \textbf{等额分配基准} & \textbf{仅评分解基准} \\
\midrule
放贷企业数 & 106 & 108 & 108 \\
投放总额/万元 & 10000 & 10000 & 10000 \\
加权平均利率/\% & 5.52 & 5.47 & 5.47 \\
利息收入/万元 & 552 & 547 & 547 \\
预期违约损失/万元 & 41 & 96 & 68 \\
\textbf{期望净收益/万元} & \textbf{511} & 451 & 479 \\
期望收益率/\% & 5.11 & 4.51 & 4.79 \\
\bottomrule
\end{tabular}
\end{table}

可以看出，本文最优策略通过"剔除高违约企业、向低风险企业倾斜额度"显著降低了预期违约损失，使期望净收益达到 511 万元，较等额分配基准提升约 13.3\%，较仅按评分分配的基准提升约 6.7\%。若进一步考虑贷款期限差异化带来的复利与展期收益，实际收益还可更高。

最终，本文给出 123 家企业的完整信贷决策清单（含是否放贷、额度、利率、期限），其中：
\begin{itemize}[itemsep=1pt]
    \item 放贷企业 106 家，合计额度 10000 万元，利率区间 $[4.42\%, 7.85\%]$；
    \item 拒绝放贷企业 17 家，主要为 D 级、存在违约记录或传染后风险超阈值的企业。
\end{itemize}

%===============================================================
% 九、灵敏度分析
%===============================================================
\section{模型灵敏度分析}

为检验模型的稳健性，本文对关键参数进行灵敏度分析。

\subsection{融合权重 $\alpha$ 的灵敏度}

综合信用得分 $S_i=\alpha C_i+(1-\alpha)(1-p_i)$ 中，$\alpha$ 反映"经营实力"与"违约概率"的相对重要性。令 $\alpha\in\{0.3,0.4,0.5,0.6,0.7\}$，重新计算信贷策略与收益，结果见表 \ref{tab:sens1}。可见在 $\alpha$ 的合理区间内，期望收益变化不超过 3\%，说明模型对 $\alpha$ 不敏感。

\begin{table}[H]
\centering
\caption{融合权重 $\alpha$ 的灵敏度}
\label{tab:sens1}
\small
\begin{tabular}{lccccc}
\toprule
$\alpha$ & 0.3 & 0.4 & 0.5 & 0.6 & 0.7 \\
\midrule
放贷企业数 & 104 & 105 & 106 & 107 & 108 \\
期望净收益/万元 & 497 & 505 & 511 & 508 & 496 \\
\bottomrule
\end{tabular}
\end{table}

\subsection{违约损失率 $\mathrm{LGD}$ 的灵敏度}

$\mathrm{LGD}$ 直接影响违约损失测算。取 $\mathrm{LGD}\in\{0.5,0.6,0.7,0.8\}$，结果见表 \ref{tab:sens2}。随着 $\mathrm{LGD}$ 上升，银行趋于收紧放贷（放贷企业数下降、利率上浮），但总体收益仍为正，模型行为符合经济直觉。

\begin{table}[H]
\centering
\caption{违约损失率 $\mathrm{LGD}$ 的灵敏度}
\label{tab:sens2}
\small
\begin{tabular}{lcccc}
\toprule
$\mathrm{LGD}$ & 0.5 & 0.6 & 0.7 & 0.8 \\
\midrule
放贷企业数 & 108 & 106 & 103 & 100 \\
加权平均利率/\% & 5.48 & 5.52 & 5.60 & 5.71 \\
期望净收益/万元 & 528 & 511 & 482 & 449 \\
\bottomrule
\end{tabular}
\end{table}

\subsection{传染强度系数 $\lambda$ 的灵敏度}

传染强度 $\lambda$ 影响问题二的策略调整幅度。$\lambda$ 越大，传染后风险抬升越明显，高中心性企业额度下调越多，但放贷集合基本稳定，说明模型结构稳健。

%===============================================================
% 十、模型评价与推广
%===============================================================
\section{模型评价与推广}

\subsection{模型优点}

\begin{enumerate}[label=(\arabic*),itemsep=2pt]
    \item \textbf{数据驱动、客观赋权}：熵权法避免了主观权重的人为偏差，TOPSIS 综合评价兼顾多维度指标。
    \item \textbf{风险与实力双视角}：融合 TOPSIS 贴近度（实力）与 Logistic 违约概率（风险），评价更全面，且 Logistic 模型可解释性强。
    \item \textbf{引入复杂网络刻画风险传染}：将发票往来建模为供应链网络，用 SIS 思想量化风险传染，弥补了孤立评价的不足，契合中小微企业"一荣俱荣、一损俱损"的生态特征。
    \item \textbf{优化模型可落地}：期望收益最大化模型目标明确、约束清晰，可直接嵌入银行信贷审批与定价系统。
\end{enumerate}

\subsection{模型缺点}

\begin{enumerate}[label=(\arabic*),itemsep=2pt]
    \item 违约样本仅 11 家，Logistic 回归存在一定的小样本偏差，违约概率估计的方差较大。
    \item 假设违约损失率、传染强度等为常数，未考虑行业、地域、时变的异质性。
    \item 未纳入宏观经济周期、政策冲击等系统性风险因素。
\end{enumerate}

\subsection{模型改进与推广}

\begin{enumerate}[label=(\arabic*),itemsep=2pt]
    \item 引入机器学习模型（随机森林、XGBoost、LightGBM）提升违约预测精度，并通过交叉验证缓解小样本问题。
    \item 将 $\mathrm{LGD}$、违约概率与行业景气度、宏观经济变量关联，构建动态风险模型。
    \item 用 SIR、SIER 等更丰富的传播动力学模型，或引入多层网络刻画不同贸易类型的差异化传染。
    \item 考虑多期动态规划与贷款组合的信用风险（如 KMV、CreditMetrics），实现风险与收益的全局最优。
\end{enumerate}

该模型可推广至供应链金融、应收账款融资、企业征信评分、地方政府债务风险传染评估等场景，具有较强的普适性。

%===============================================================
% 参考文献
%===============================================================
\begin{thebibliography}{9}
\bibitem{ref1} 姜启源, 谢金星, 叶俊. 数学模型（第五版）[M]. 北京: 高等教育出版社, 2018.
\bibitem{ref2} 司守奎, 孙玺菁. 数学建模算法与应用（第2版）[M]. 北京: 国防工业出版社, 2015.
\bibitem{ref3} Hwang C L, Yoon K. Multiple Attribute Decision Making: Methods and Applications[M]. Berlin: Springer, 1981.
\bibitem{ref4} Shannon C E. A mathematical theory of communication[J]. Bell System Technical Journal, 1948, 27(3): 379--423.
\bibitem{ref5} Hosmer D W, Lemeshow S. Applied Logistic Regression[M]. New York: Wiley, 2000.
\bibitem{ref6} Brin S, Page L. The anatomy of a large-scale hypertextual Web search engine[J]. Computer Networks and ISDN Systems, 1998, 30(1--7): 107--117.
\bibitem{ref7} Pastor-Satorras R, Castellano C, Van Mieghem P, et al. Epidemic processes in complex networks[J]. Reviews of Modern Physics, 2015, 87(3): 925--979.
\bibitem{ref8} 杨晓光, 李自然. 供应链金融风险管理研究综述[J]. 管理评论, 2015, 27(5): 135--144.
\bibitem{ref9} 全国大学生数学建模竞赛组委会. 2020 年高教社杯全国大学生数学建模竞赛 C 题[EB/OL]. http://www.mcm.edu.cn.
\end{thebibliography}

%===============================================================
% 附录
%===============================================================
\appendix
\section{核心代码（Python）}

\subsection{数据清洗与指标提取}

\begin{verbatim}
import pandas as pd
import numpy as np

# 读取发票与企业信息
invoice = pd.read_excel('附件1_进销项发票.xlsx')
firm = pd.read_excel('附件2_企业信息.xlsx')

# 清洗：剔除作废发票
invoice = invoice[invoice['发票状态'] != '作废']

# 价税合计
invoice['交易额'] = invoice['金额'] + invoice['税额']

# 分离销项与进项
sale = invoice[invoice['方向'] == '销项']
buy  = invoice[invoice['方向'] == '进项']

# 企业实力指标
def build_index(firm_id):
    s = sale[sale['企业代号'] == firm_id]
    b = buy[buy['企业代号'] == firm_id]
    X1 = s['交易额'].sum()              # 净销售额
    X2 = b['交易额'].sum()              # 净采购额
    X3 = len(s)                          # 销售笔数
    X4 = len(b)                          # 采购笔数
    X5 = invoice['月份'].nunique()       # 活跃月份数
    X6 = pd.concat([s['对手'], b['对手']]).nunique()  # 交易对手数
    # 下游集中度 HHI
    share = s.groupby('对手')['交易额'].sum() / X1
    X7 = (share ** 2).sum()
    X8 = X1 / X2 if X2 != 0 else 0       # 销采比
    X9 = (s['交易额'] < 0).sum() / X3    # 作废/退货率
    return [X1, X2, X3, X4, X5, X6, X7, X8, X9]
\end{verbatim}

\subsection{熵权法--TOPSIS}

\begin{verbatim}
def entropy_weight(X):
    X = X / X.sum(axis=0)          # 比重
    e = -(X * np.log(X + 1e-12)).sum(axis=0) / np.log(len(X))
    d = 1 - e
    return d / d.sum()              # 熵权

def topsis(X, w):
    Z = X * w
    Zp, Zn = Z.max(axis=0), Z.min(axis=0)
    dp = np.sqrt(((Z - Zp) ** 2).sum(axis=1))
    dn = np.sqrt(((Z - Zn) ** 2).sum(axis=1))
    return dn / (dp + dn)           # 相对贴近度
\end{verbatim}

\subsection{Logistic 回归违约概率}

\begin{verbatim}
from sklearn.linear_model import LogisticRegression

# 特征：评级数值化 + 贴近度 + 对数销售额 + 对手数
X_feat = df[['评级数值', '贴近度', 'log销售额', '对手数']]
y = df['是否违约'].astype(int)

model = LogisticRegression(C=1.0)
model.fit(X_feat, y)
p_i = model.predict_proba(X_feat)[:, 1]   # 违约概率
\end{verbatim}

\subsection{供应链风险传染}

\begin{verbatim}
def risk_contagion(R0, W, lam=0.35, tol=1e-6, max_iter=100):
    R = R0.copy()
    for _ in range(max_iter):
        Rnew = R0.copy()
        for i in range(len(R0)):
            # 上游 j 传染给 i
            Rnew[i] += lam * sum((R[j] - R0[j]) * W[j, i]
                                 for j in range(len(R0)))
        if np.max(np.abs(Rnew - R)) < tol:
            break
        R = Rnew
    return R
\end{verbatim}

\subsection{遗传算法求最优信贷决策}

\begin{verbatim}
import random

def fitness(genes, p, S, Lmax=10000, rf=0.0435, beta=0.06, LGD=0.6):
    x = np.round(genes[:, 0]); L = genes[:, 1]
    L[x < 0.5] = 0
    r = rf + beta * (1 - S)                 # 风险溢价定价
    profit = L * (r * (1 - p) - p * LGD)    # 单户期望利润
    total = L.sum()
    pen = max(0, total - Lmax) * 100        # 超约束惩罚
    return (profit * (x >= 0.5)).sum() - pen

def ga(p, S, pop=200, gen=100):
    genes = [np.random.rand(len(p), 2) * 500 for _ in range(pop)]
    for _ in range(gen):
        fit = [fitness(g, p, S) for g in genes]
        # 轮盘赌选择、单点交叉、高斯变异（略）
        ...
    return genes[np.argmax(fit)]
\end{verbatim}

\end{document}
